% PREFACE -- the multi-query surrogate problem, scope, intended reader, part
% dependencies, two reading maps, notation conventions, deliberate omissions.
% STATUS: authored; adversarially verified against the .toc and every chapter, section,
% appendix, definition, assumption and table number named below. Four false claims fixed:
% prerequisites sit in Ch. 1's opener, not Sec. 1.1; Ch. 16 does develop one calibration;
% Ch. 22's baseline is FBP; and Ch. 10 assumes CUDA/cuDNN, so no toolchain omission.
\chapter*{Preface}
\addcontentsline{toc}{chapter}{Preface}
\markboth{Preface}{Preface}

Numerical simulation answers one question at a time: data go in, the matching solution field
comes out, and nothing carries over to the next run. The workloads that organize this book do
not pose one question. Shape optimization queries the same equations for thousands of
geometries, ensemble forecasting integrates fifty perturbed initial states twice a day, and
posterior sampling needs a forward evaluation at every proposal a Markov chain makes.
Simulation costs $M\,c_{\mathrm{sim}}(N)$, a surrogate $c_{\mathrm{train}} +
M\,c_{\mathrm{eval}}(N)$: $M$ queries, one classical solve at $N$ unknowns, one training run,
one evaluation per query, as Section~1.1 sets out in \eqref{eq:ch01-cost-bookkeeping}. I wrote
this book because the ledger has a catch: the saving is real only if the surrogate can be asked
the question at all, and a network whose weights are sized for $N$ grid points has no defined
action on $N'$. What must be learned is therefore not a solution but the map from data to
solution --- a map on function spaces, stated without reference to a mesh, evaluated on
whichever discretization a query supplies. Amortization is what changes; the resolution
dependence remains.

Part~I (Chapters~1--9) builds the framework that makes such a map learnable: the Sobolev scale
and compact embeddings that make uniform approximation meaningful in infinite dimensions
(Chapter~2); the proof that the Darcy, Burgers and forced Navier--Stokes solution operators are
continuous, stable and smoothing (Chapter~3); the density theorems of Chen \& Chen and their
DeepONet and generic-operator counterparts (Chapter~4); spectral truncation bounds and
convergence rates for the Fourier neural operator, discretization invariance and the continuum
limit, and physics-informed and Sobolev-regularized objectives (Chapters~5--7); a taxonomy
of what each architecture requires of its target and where its theorems are (Chapter~8), and
the open questions (Chapter~9). Part~II (Chapters~10--16)
turns those objects into tensor programs: layouts, coordinate encoding, losses and normalization
(Chapter~10); complete implementations of DeepONet, the FNO, and the graph and geometry-informed
operators (Chapters~11--13); the canonical benchmarks with their specifications (Chapter~14);
the regimes and failure modes a long run meets (Chapter~15); and export, cost accounting and
calibration (Chapter~16). Part~III (Chapters~17--24) leaves the periodic box for triangulated
surfaces, the sphere, faulted reservoirs, patient-specific ventricles and undersampled
sinograms, scored against the solvers they displace --- finite element, finite volume, implicit
finite difference, spectral-transform, and, in Chapter~22, filtered backprojection --- with
boundary enforcement as the question they make unavoidable (Chapter~23) and a synthesis matrix
closing the argument (Chapter~24). The appendices hold what the text borrows rather than proves,
from the functional-analysis reference (Appendix~A) to the glossary (Appendix~F).

The intended reader is a graduate student or researcher in scientific machine learning who wants
to read a convergence proof and a training log without a translation step. The prerequisites are
those Chapter~1 states before its first section, and they are concrete. From analysis: the
measure-based functional analysis Section~1.2 uses without apology --- a Borel probability
measure on the input space, an expectation of a relative error functional, a worst-case bound
over a compact set held distinct from a risk under a measure --- plus Chapter~2's machinery:
Riesz representation, duality pairings, weak derivatives, Sobolev spaces and their embeddings at
the level of \cite{brezis2011functional, evans2010partial}. From numerical analysis: the
discrete Fourier transform as a change of basis --- what a mode index means, what Nyquist
forbids. From practice: enough PyTorch to read a layer definition --- a \texttt{torch.nn.Module}
and a batched tensor contraction \cite{paszke2019pytorch}, which is also what Chapter~10 assumes.
Readers who command the analysis may treat Chapter~2 as review, as Chapter~1's close suggests;
the rest should not skip it: Part~I states every theorem in that language.

The parts are not independent. Part~II takes Part~I's discretization-invariance framing as its
design criterion: Chapter~10 opens on Definition~6.1 and treats layout, coordinate encoding, loss
and normalization as what a program owes that definition; Chapter~16 asks whether an exported
graph still computes its sizes from the tensor handed to it. The dependence runs the
other way at the end of Part~I, where Chapter~9 turns its unsolved problems into a
specification: the truncation tail becomes a choice of \texttt{modes}, the unrolled loss a
training recipe, the Lipschitz product a divergence the practitioner must notice by hand.
Part~III uses both. Chapter~17 names the concessions of Chapters~1--16 --- torus, unit square,
one application of the operator --- and withdraws them at mesh sizes two to five orders above
those Chapter~16 counted. Each application chapter needs Part~II for the artifact and Part~I for
the hypothesis it cannot escape: the bounds of Chapters~4--6 are stated under a compactness
hypothesis on the support of the training measure (Assumption~9.1), so out-of-distribution
behavior is recorded in Part~III as observed, not bounded.

Two routes are worth prefiguring. For a one-semester course, Chapters~1--6 form the spine; add
Chapters~10 and~12 for the implementation, Appendix~B as the laboratory exercise --- its $0.0141$
relative $L^{2}$ errors on the $64$- and $256$-point grids, caveat attached --- and Chapter~14
with Appendix~C for the specifications a run needs. One application chapter, read against
Chapter~17's ledger, closes the course; Chapter~19 or~20 suits best. Where time is short,
Chapter~8's theorem ledger is the hour best spent and Chapter~9 the discussion. For a
practitioner with a model to train, begin at Chapter~10; take Chapter~12 (Section~12.4 for the
official library), Chapter~11 if a branch-and-trunk fit is intended, Chapter~13 for meshes, then
Chapter~15 with Appendix~D before the first long run and Chapter~16 before serving. Read two
theory chapters regardless: Chapter~6 for what a discretization owes the continuum, Chapter~5
for what the choice of modes costs. Chapter~4's three-term split explains a validation curve
that flattens while the training loss still falls; Chapter~9, why a rollout that looks stable is
not bounded.

Notation is fixed once, in \emph{Notation and Conventions} and its Table~\ref{tab:notation-master},
which gives the symbols in three blocks: Part~I's infinite-dimensional objects, Part~II's tensor
layouts, Part~III's physical quantities. Its rule is that a letter reused across parts carries
its ambiguity in its own row rather than silently: $\mu$ is the input measure, Chapter~10's
pointwise ensemble mean, and, with a phase subscript, a Part~III viscosity; $M$ counts queries in
Chapter~1, ensemble members in Part~II, detector views in Chapter~22. Where a chapter's gloss and
the table differ, the table governs. Appendix~E inverts it: Section~E.1 indexes the variants (GNO,
GINO, AFNO, U-FNO, PINO) by the chapter that builds each, Section~E.2 maps the
principal symbols to the equation or section at which they begin to do work.

What the book does not cover deserves the same directness. There is no learning theory beyond
Chapter~9's entropy arithmetic: Chapter~4's density theorems are existence results whose limits
Section~4.4 inventories, and optimization error, Section~1.3's fourth term, is
handled empirically in Chapter~15 and Appendix~D alone. Supervision is by paired fields; the one
route around it is PINO's weak-form residual objective (Chapter~7, Section~14.4), whose reported
gains are benchmark outcomes, not theorems. Sparse or partially observed labels are not treated:
DeepONet's sensors (Section~11.2) restrict where an input is read, not whether an output is
known, and Chapter~22's sparse-view and limited-angle protocols thin the sinogram, not the
labels. Inverse problems appear as workloads and as pointers (Sections~1.2, 3.1, 17.2 and~20.4;
Chapter~9's fallback to Bayesian formulations), since this book learns forward maps, not
posteriors; the one inverse problem it solves, Radon inversion in Chapter~22, is linear.
Uncertainty quantification is likewise a use case (Chapter~20's Monte Carlo sweeps) and a
roadblock (Chapter~24's bound on distribution shift), not a methodology: no posterior is
estimated, no coverage guarantee proved, and the only calibration developed is
Chapter~16's, of surrogate error against its labeling budget. Classical solvers
are baselines and label producers; Section~23.1's account of why a conforming
discretization carries its boundary conditions by construction is the one place their mechanics
matter; elsewhere they are priced rather than derived: Chapter~17 charges meshing as an
unreported cost, the turbulence and capillary-pressure closures only as far as the surrogate's
data depend on them. Software engineering stops where Chapter~16 stops: no cluster scheduling,
data pipeline or model-versioning practice, and no tuning of a vendor's kernels beyond the CUDA
and cuDNN vocabulary Chapter~10 assumes. Finally, every surrogate here is deterministic: no
generative model --- diffusion, flow, or sampler-based --- is treated, and \emph{diffusion} in
Part~III names a physical coefficient.

Three conventions govern the text, which is why those omissions are named: numbers travel with
the configuration that produced them, theorems with their hypotheses, and a claim whose
provenance cannot be recovered is recorded as unestablished rather than smoothed. Where the
record is silent --- a baseline timing behind a published speedup --- the book says so.
